\subsection{19.61: root closure preserves the derived carpet}
\label{cand:19.61}
\problemstatement{19.61}

Use the notation of Section~\ref{cand:19.62}. For a carpet \(A\), put
\[
 H=\langle x_r(A_r):r\in\Phi\rangle,\qquad
 D_r=\{t\in K:x_r(t)\in H\}.
\]
The candidate asserts that \(D\) is a carpet and
\(\partial D=\partial A\), over every commutative ring.
Troyanskaya previously proved carpet closedness of this root
closure in odd characteristic, with characteristic divisible by
three also excluded in type \(G_2\) \cite{troyanskaya}.
The argument here uses Nuzhin's known square criterion and a
universally valid finite adjoining lemma. Priority of the stronger
statement is not established by those source comparisons.

\begin{lemma}
Let \(B=\partial A\), fix \(a\in A_r,b\in A_{-r}\), and replace
\(A_r\) by \(A_r+\Z a^2b\), leaving all other entries unchanged.
The resulting family \(A'\) is a carpet and \(\partial A'=B\).
\end{lemma}
\begin{proof}
Consider an ordered Chevalley rule whose first input changes from
\(x\) to \(x+na^2b\). Its parameter expands as
\[
 c(x+na^2b)^iu^j
   =\sum_{k=0}^i c\binom ik n^k x^{i-k}a^{2k}b^ku^j.
\]
The \(k=0\) term lies in \(B_{ir+js}\) by definition. For every
other term, the required universal implication is
\begin{equation}
 c\binom ik x^{i-k}a^{2k}b^ku^j\in B_{ir+js},
 \quad x,a\in A_r,\ b\in A_{-r},\ u\in A_s.
 \label{eq:adjoin-target}
\end{equation}
The variables \(x,a\) are independent even though their root
labels coincide. Multiplication by any integer \(n^k\) is allowed
in the additive subgroup \(B_{ir+js}\).
The analogous substitution into the second input is checked
separately, retaining the constant of the original ordered rule;
reversing an ordered source pair in \(G_2\) can change its absolute
constant. A nontrivial rule has distinct nonopposite source roots,
so only one input subgroup changes.

The same root-plane argument as in Section~\ref{cand:19.62}
reduces \eqref{eq:adjoin-target} to types \(A_2,B_2,G_2\).
The ancillary certificates apply exactly the integer proof
calculus explained there. Their complete sizes are
\[
\begin{array}{c|rrr}
 \text{type}&\text{targets}&\text{nodes}&\text{largest derivation}\\ \hline
 A_2&12&72&6\\
 B_2&48&380&12\\
 G_2&240&2739&19\\ \hline
 \text{total}&300&3191&19
\end{array}
\]
Every noninitial node lies in \(B\) by a carpet application and,
where needed, a Bezout combination of two such applications.
The last coefficient divides the coefficient in
\eqref{eq:adjoin-target}. Thus each certificate is a finite
integer algebraic proof valid under every substitution in \(K\).
They prove all 300 implications. For example, in type \(A_2\)
with \(p=r+s\), the only changed first-input term is \(a^2bu\).
The rules give \(au\in B_p\), then \(abu\in B_s\), and then
\(a^2bu\in B_p\). The longer certificates perform the same
operation where coefficients two and three intervene.

The coverage checker reconstructs the requests by repeated sparse
polynomial multiplication, independently of the binomial formula
used by the target generator. It covers both inputs of all ordered
rules: 416 substitutions and 600 nonconstant terms, giving the
300 requests after coincident requests are combined by gcd.
The checked rule tables and the complete derivations accompany
this paper; their search procedure is not a mathematical premise.

Expansion now gives \(\partial A'\subseteq B\); monotonicity gives
the reverse inclusion. Since \(B\subseteq A\subseteq A'\), every
defining parameter of \(A'\) lies in its designated entry.
This is precisely the carpet condition.
\end{proof}

Starting with \(C^{(0)}=A\), define
\[
 C^{(n+1)}_r=C^{(n)}_r+
 \langle u^2v:u\in C^{(n)}_r,\ v\in C^{(n)}_{-r}\rangle_{\Z}.
\]
Each \(C^{(n)}\) is a carpet with derived family \(B\).
Indeed, adjoin any finite collection of the displayed generators
one at a time. Its parameters remain available after earlier
adjoinings, so the lemma applies throughout. These finite
extensions form a directed family whose union is \(C^{(n+1)}\).
A carpet condition involves finitely many elements, and a derived
parameter uses only two input elements. Hence directed union
preserves both carpet membership and equality of the derived family.

Put \(C_r=\bigcup_n C^{(n)}_r\). Again \(C\) is a carpet and
\(\partial C=B\). Any \(u\in C_r,v\in C_{-r}\) lie in a common
stage, so \(u^2v\) lies in the next stage. Thus \(C_r^2C_{-r}\subseteq
C_r\), and the known square criterion makes \(C\) closed.
Since \(A\subseteq C\), one has \(H\subseteq E(\Phi,C)\);
closedness yields
\[
 A\subseteq D\subseteq C.
\]
The root-group addition law makes each \(D_r\) additive.
Every commutator parameter formed from \(D\) belongs to
\(\partial C=B\subseteq A\subseteq D\), proving that \(D\) is
a carpet. Finally monotonicity gives
\[
 B=\partial A\subseteq\partial D\subseteq\partial C=B.
\]
The constructed square completion \(C\) need not equal \(D\);
the proof only requires the displayed sandwich. It does not assert
that the adjoined root elements already belong to \(H\).
In rank one there are no nonopposite commutator conditions and
all derived families are zero, so the conclusion is immediate.
See \artifact{research/19.61-proof.md} and
\artifact{results/19.61-square-completion-packet.json}.
